\section*{Discussion}

\subsection*{What the evaluation supports}
The deployment produced many serious errors, and its own audit caught few
of them. The written sweep report of the in-session audit fully detected
10\% of the 68 reference errors and none of the 6 critical errors among
them. Across the full ledger, the report flagged none of the 11 critical
errors that were already present when the audit ran (one of them only
partly present). Later, during the audit's second repair round, the
deployment agent found one of these 11: the direction-of-change metric.
The other 2 of the 13 critical errors were made during the
audit's repair rounds. Fresh agents that audited the same snapshot found four times as
many reference errors, whether they used the checklist or a generic
instruction to review critically.

In the pre-registered comparisons with Claude Opus~5.5, the checklist did
not clearly help reviewers find more problems than a generic review of
the same budget in any of the three studies, and in none did it reduce
false alarms. There were two partial exceptions. In the exploratory arm
with Claude Haiku~4.5, a checklist review followed by repair raised the
share of pitfalls handled by 6.4 percentage points more than a generic
review followed by repair (95\% bootstrap interval over the 28 paired
executor runs, pooled over tasks, 1.3 to 11.7), but it did not raise the
share of correct verdicts. In Study~C the two checklist arms (generic and
as deployed) fully detected more of the 17 missing baselines and nulls
(51\% and 55\% against 39\% for the generic review) and of the 9 wrong
statistics (22\% and 37\% against 15\%), but fewer of the 12 code bugs
(69\% and 58\% against 81\%); these splits by error type were not tested
for significance. The specifications did not change whether a current
agent reached the correct conclusion on focused analysis tasks: with
Claude Opus~5.5 every executor run reached the correct verdict with or
without them. With a smaller model, Claude Haiku~4.5, they may help:
correct verdicts rose from 10 to 15 of 20 runs, a difference that 20 runs
per arm cannot confirm. The framework as deployed therefore did not make
the agent's work reliable, and the checklist added little over a generic
review of the same budget.

What did catch errors was independent review in a fresh context. The
blind re-audits of Study~C and the verification that built the ledger
were carried out by agents that had not produced the work and had no
stake in it. The in-session audit, by contrast, was written and applied
by the deployment agent itself, in one working session on the last day of
the deployment. Apart from a same-day extension of the topological
hypothesis screen, the work it audited had been done in other sessions.
The audit's sweep of all eight pipelines took about
8~minutes from request to report (3.7~hours counting the two repair
rounds), against 21 to 29~minutes for each blind re-audit, and it graded
run summaries more than code. We cannot separate the effects of
independence, more time and a different model: the deployment and its
in-session audit ran on Claude Opus~4.7, and the re-audits on Claude
Opus~5.5. Whatever the cause, an
audit by the agent that did the work was weak evidence that the work was
correct.

\subsection*{Most of the errors that mattered were in code}
Seven of the 13 critical errors were code bugs: an intervention hook that
edited the wrong tensor and deleted a transformer block (three ledger
entries), combined edits that overwrote each other, a model input encoded
on the wrong scale, a metric that ignored the prediction, and a lookup
that read the wrong embedding rows. None of these is a methodological
pattern of the kind the checklist describes, and the checklist did not
help fresh reviewers find them (the generic review found more of the code
bugs). Each of the seven, however, would have been caught by a short
deterministic test written for the analysis: an edit of zero must change
nothing; an edit at layer~$l$ must leave earlier layers unchanged;
combined edits must differ from single edits; the token order must match
the model's reference encoder on a few cells; a direction-of-change
accuracy must be compared with the constant-sign predictor on the same
records; and for a few named genes, the embedding row read must be the
row that the model's own vocabulary file gives them. Such tests are cheap
and do not depend on an agent's judgment. The other six critical errors
were not code bugs: three were missing or wrong nulls or baselines, two
were wrong statistics (one of them in the audit's own summary) and one
was an invalid proxy axis. The ledger places five of these six under a
checklist pattern (P1, P3 or P5 in Table~\ref{tab:checklist}; S1~Table).
The generic deterministic checks we built (structure of the
specification, parameter deviations, number tracing) catch a different
class of problem. By design they do not test whether a null randomises,
an edit lands on the right tensor, an input is encoded as the model
expects, a lookup reads the right rows or a number means what the text
says, so they would not have caught 12 of the 13 critical errors. The thirteenth, the audit's percentage with
no cell-by-cell record behind it, is the kind of problem the number
tracer looks for. We ran the tracer only on the run summaries, not on the
audit reports, so we did not test whether it would have flagged it. For
agent-run interpretability we would therefore recommend
analysis-specific tests of interventions, inputs, lookups and baselines,
run before results are summarised, together with independent
fresh-context review, rather than a longer checklist.

The held-out benchmark adds a second observation. Every review of a
flawed, single-analysis package found its documented error. This
included a raw-count input fed to an encoder that expects binned values,
the same kind of error as MaxToki's input encoding. Only part of
the MaxToki input-encoding error was in the Study~C reference set: three attention runs
fed $\log(1+\text{CP10k})$ values in as if they were counts (S1~Table,
L118). Of the nine blind audits of the whole MaxToki project, three were
graded as fully detecting this part and two as detecting it partly. These
five suspected that the Adamson input was log-normalised and noted that
no code checked for raw counts. None found that every run except the RPE1
attention run was encoded on the wrong scale.
Reviewing one analysis at a time, with its code, outputs and summary
together, appears far more effective than reviewing a large project in
one pass. This comparison is across studies, with different budgets per
analysis, and should be tested directly.

Reviews also over-called errors. Most reviews of correct analyses (53 of
66) claimed at least one error in a main result that was not one. On
correct and flawed analyses together, the reviews made 105 such claims.
The graders marked 82 of them as false alarms (73 on correct and 9 on
flawed analyses); the verifier agent ruled the other 23 false alarms
after a grader dispute. A post-hoc check of the 82, which was not
pre-registered, found that 76 were accurate observations whose severity
was overstated. A review step
is therefore only as good as the step that checks its points against the
data before acting on them. In Study~A a repair step was meant to do
this, and it did not always succeed. With Opus~5.5, repairs after generic
reviews turned 3 of 40 correct verdicts into wrong ones, and repairs
after checklist reviews turned none. In all three cases the reviewer
leaned on metrics named in the specification, one of which has a high
chance level that the specification does not state. A specification can
thus mislead a reviewer even when the executor that followed it was not
misled. In the exploratory Haiku~4.5 arm the pattern was reversed:
repairs after checklist reviews turned 3 of 21 correct verdicts wrong
(and 3 wrong verdicts right), and repairs after generic reviews turned
none wrong. So each review style, with one of the two models, led
repairs to turn correct verdicts wrong.

\subsection*{Specifications and prompts}
During the deployment a specification was, in practice, a long and
structured prompt: nothing checked that it was complete, that its
parameters were used or that its validation criteria were met, and the
executing agent changed parameters (for example, the number of null
draws) without recording why in most cases. What can be checked by code
is narrow (Table~\ref{tab:checks}): document structure, resolvable code
references, deviations from machine-readable parameters, and whether a
stored gate flag matches the stored numbers. Whether a null randomises,
an intervention edits the intended tensor or a number means what the text
says remains a matter of judgment unless an analysis-specific test is
written for it. The specifications also carry expectations from earlier
studies; one stated that a negative verdict was the expected outcome. For
a strong agent on a focused task this made no difference to the verdict,
but a contract that states the expected answer is a source of bias that a
pre-registered analysis plan would avoid.

\subsection*{Transfer to other models and settings}
The held-out items came from other projects by the author: a text-model
cell embedding (Cell2Sentence), a protein language model (ESM-2), published
single-cell transformers (scGPT, Geneformer and STATE), and
gene-regulatory-network benchmarks that used no model. Every documented
error was found in each of these four groups. In three groups most
reviews of correct analyses raised a false alarm: 18 of 18 reviews for
Cell2Sentence, 12 of 18 for ESM-2 and 17 of 18 for single-cell
transformers. In the gene-regulatory-network group half did (6 of 12).
Each group had only two or three correct items, too few to compare the
groups. Because the checklist's patterns concern study
design rather than any model's internals, we expect the review findings
to transfer to other foundation models in biology. The code-level tests do not transfer
unchanged: each architecture and tokenizer needs its own tests of where an
edit lands and how inputs are encoded. We tested one agent model family
(Claude) as executor, reviewer and grader, one deployment, and analyses
from one author's projects; other agent families, other research groups'
analyses and human reviewers were not tested.

\subsection*{No human or non-agent baseline}
We did not compare the agent workflow with human analysts or human
reviewers, and we make no claim about efficiency. The deployment's
supervision time was not logged, and the controlled studies measured
agent effort (tool calls, time), not human effort. Whether a human
reviewer with the same budget would find more of the code-level errors,
or raise fewer false alarms, is an open question that a study with human
reviewers could answer on the same items, which are released for this
purpose.

\subsection*{What the case study says about MaxToki-217M, and what it does not}
The corrected analyses describe MaxToki-217M given one cell at a time,
with rank-value inputs from K562, RPE1 and Tabula Sapiens cells, and
judged against TRRUST, DoRothEA and CRISPR-interference screens. They do
not test MaxToki's multi-cell trajectory mode, the larger MaxToki-1B
(except for one attention run, whose inputs were encoded incorrectly and
which we do not report), or other single-cell foundation models. Negative
results here mean that these probes, with these references, found no
signal beyond simple baselines in this model; they are not evidence that
\scfm{}s in general lack regulatory information. The references themselves
are limited: TRRUST covers few transcription factors per cell line, and
the knockdown fold-changes are noisy, so modest signals could be missed.

\subsection*{Limitations}
The executor, reviewer, repairer and grader agents all came from one model
family, so shared blind spots cannot be excluded; graders agreed closely
with each other, but no human graded the deliverables. Study~A showed a
ceiling for Opus~5.5, which leaves it without power to detect benefits of
the specification for that model; the exploratory Haiku arm was added
after that ceiling was seen and has 20 runs per arm (S2~Appendix). The Study~A tasks were built from the
deployment's data and so are not independent of it. The Study~B items
came from the author's own projects. Each flawed item held one documented
error, and every review found it, in both arms (33 of 33 per arm). So
Study~B could not show whether the checklist finds more of the documented
errors. It could compare only false alarms and the number of other real
problems found per review (2.1 in each arm). The error ledger is a lower bound:
blind re-audits added 60 errors to it and the full input-encoding error was
found last, so detection rates are relative to the errors known when
each study's key was frozen. The deployment ran on Claude Opus~4.7 (its
one automated loop called Claude Sonnet), while the evaluation used
Opus~5.5 and, in the exploratory arm, Haiku~4.5. So a difference in
model may explain part of the gap between the in-session audit and the
blind re-audits. Finally, this manuscript
was drafted by an agent under the author's direction. The source data of
every figure, the error ledger with the file location of each error
(S1~Table), and a results map that lists every number in the paper with
its source file (RESULTS\_MAP.csv) are released with the code.
